
\documentclass[11pt]{article}
\usepackage[a4paper,margin=1in]{geometry}
\usepackage[T1]{fontenc}
\usepackage{lmodern}
\usepackage{xcolor}
\usepackage{listings}

\lstdefinestyle{plain}{
  basicstyle=\ttfamily\small,
  columns=fullflexible,
  keepspaces=true,
  frame=single,
  showstringspaces=false,
  breaklines=true
}

\title{Type-System Errors}
\author{}
\date{}

\begin{document}
\maketitle

\section*{Type-System Errors Only}

\begin{enumerate}
  \item \textbf{Kind mismatch} --- using a palette where a graph is expected.
\begin{lstlisting}[style=plain]
transform t = toUDG(p3)
\end{lstlisting}

  \item \textbf{Vertex not in graph} --- restriction mentions a vertex outside the graph.
\begin{lstlisting}[style=plain]
graph g = { vertices=[1,2,3], edges=[(1,2)] }
restrictions r = { precolor 4 : red; }
\end{lstlisting}

  \item \textbf{Edge endpoint not in graph} --- edge uses an undeclared vertex.
\begin{lstlisting}[style=plain]
graph g = { vertices=[1,2,3], edges=[(1,4)] }
\end{lstlisting}

  \item \textbf{Color not in palette} --- restriction uses a color unavailable in the palette.
\begin{lstlisting}[style=plain]
palette p = colors [red, blue]
restrictions r = { precolor 1 : green; }
\end{lstlisting}

  \item \textbf{Palette-size contradiction} --- restriction asks for more colors than the palette provides.
\begin{lstlisting}[style=plain]
palette p = colors [red, blue]
restrictions r = { useExactly 3; }
\end{lstlisting}

  \item \textbf{Wrong graph in restriction set} --- a restriction set typed for one graph is reused on another.
\begin{lstlisting}[style=plain]
graph g1 = { vertices=[1,2,3], edges=[(1,2)] }
graph g2 = { vertices=[a,b], edges=[(a,b)] }
restrictions r1 = { precolor 1 : red; }
instance x = color g2 using p subject-to r1
\end{lstlisting}

  \item \textbf{Wrong graph in keep/drop} --- transform options mention vertices not in the input graph.
\begin{lstlisting}[style=plain]
graph g = { vertices=[1,2,3], edges=[(1,2)] }
transform t = toUDG(preserve chromatic, drop=[4])
\end{lstlisting}

  \item \textbf{Overlapping keep/drop} --- the same vertex is both kept and dropped.
\begin{lstlisting}[style=plain]
transform t = toUDG(preserve chromatic, keep=[1,2], drop=[2,3])
\end{lstlisting}

  \item \textbf{Exact transform with elimination} --- exact mode cannot drop vertices.
\begin{lstlisting}[style=plain]
transform t = toUDG(exact, drop=[4,5])
\end{lstlisting}

  \item \textbf{Guarantee mismatch} --- using a weak transform as if it preserves a stronger property.
\begin{lstlisting}[style=plain]
transform t = toUDG(preserve kColorable(3))
-- later treated as preserve chromatic
\end{lstlisting}

  \item \textbf{Wrong value in vertex position} --- a color is used where a vertex is required.
\begin{lstlisting}[style=plain]
restrictions r = { precolor red : blue; }
\end{lstlisting}

  \item \textbf{Ill-typed UDG literal} --- points/edges do not match the UDG vertex set.
\begin{lstlisting}[style=plain]
udg {
  radius = 1.0,
  points = [1@(0,0), 4@(1,0)],
  edges  = [(1,2)]
}
\end{lstlisting}
\end{enumerate}

\end{document}